\section{Introduction}
\label{sec:intro}

% 第一段：問題與重要性。
Serving large language models (LLMs) is now one of the largest consumers of GPU time.
Each request runs in two phases with different bottlenecks.
The \emph{prefill} phase processes the whole prompt at once and is compute bound.
The \emph{decode} phase then produces one token per step; each step reads all model
weights and the key--value (KV) cache, so it is bound by memory bandwidth.
Engines such as \vllm~\cite{kwon2023pagedattention} keep GPUs busy with continuous
batching~\cite{yu2022orca} and paged KV-cache management, and they split a model across
GPUs with Megatron-style tensor parallelism (TP)~\cite{shoeybi2019megatron}.

% 第二段：現有方法的不足。
TP divides every layer's matrix multiplications across $p$ GPUs, so each GPU reads only
$1/p$ of the weights.
For memory-bound decode this promises close to a $p$-fold reduction in step time.
The price is two all-reduce operations per Transformer layer, on the critical path of
every step.
On datacenter GPUs, NVLink makes this cost small.
Commodity GPUs, however, communicate over PCIe and often without direct peer-to-peer
access, so the all-reduce can erase the gain.
Prior analytical work on TP inference targets TPU pods~\cite{pope2023efficiently}, and
serving-system papers usually report results on a single interconnect.
A practitioner who owns a few consumer GPUs therefore cannot easily tell whether TP will
help, or by how much.

% 第三段：本文方法概述。
In this paper we answer two questions:
\textbf{(Q1)} How well does TP in \vllm scale on NVLink-connected H100 GPUs compared with
PCIe-connected RTX~4090 GPUs, and how much of the gap is due to communication?
\textbf{(Q2)} Can a simple model, calibrated on one GPU plus an all-reduce
microbenchmark, predict multi-GPU performance?
We run the same model, the same \vllm version and the same workload on both platforms,
sweep the batch size and the TP degree, and profile the GPU kernels.

% 貢獻列表（PPoPP 要求明確寫出 key contributions）
This paper makes the following contributions:
\begin{itemize}
  \item A controlled measurement of TP strong and weak scaling in \vllm on NVLink (H100,
        $p\le2$) and PCIe (RTX~4090, $p\le4$), separated into prefill and decode
        (\S\ref{sec:eval}).
  \item A kernel-level breakdown and an ablation of \vllm's custom all-reduce and CUDA
        graphs, which attribute the scaling gap to communication and synchronization
        (\S\ref{sec:eval}).
  \item A communication-aware analytical model that predicts TP step time from
        single-GPU runs and an all-reduce microbenchmark, with \tbd\% mean error
        (\S\ref{sec:design}, \S\ref{sec:eval}).
  \item An open benchmark harness that runs within strict shared-cluster limits
        (6-minute jobs, 2 GPUs, 100\,GB storage) and regenerates every figure in this
        paper with one command.
\end{itemize}

% 雙盲提醒：引用自己的舊作請用第三人稱。
